\documentclass[12pt, a4paper]{article}

% ─── Encoding e lingua ────────────────────────────────────────────────────────
\usepackage[T1]{fontenc}
\usepackage[utf8]{inputenc}
\usepackage[italian]{babel}

% ─── Layout ───────────────────────────────────────────────────────────────────
\usepackage[top=2.5cm, bottom=2.5cm, left=2.8cm, right=2.8cm, headheight=15pt]{geometry}
\usepackage{parskip}
\setlength{\parindent}{0pt}

% ─── Font ─────────────────────────────────────────────────────────────────────
\usepackage{lmodern}

% ─── Grafica e colori ─────────────────────────────────────────────────────────
\usepackage{graphicx}
\usepackage[table]{xcolor}
\definecolor{headerbg}{RGB}{30, 80, 140}
\definecolor{rowalt}{RGB}{235, 242, 252}
\definecolor{best}{RGB}{198, 239, 206}
\definecolor{tied}{RGB}{255, 235, 150}

% ─── Tabelle ──────────────────────────────────────────────────────────────────
\usepackage{booktabs}
\usepackage{array}
\newcolumntype{C}[1]{>{\centering\arraybackslash}p{#1}}
\newcolumntype{L}[1]{>{\raggedright\arraybackslash}p{#1}}

% ─── Link e metadati ──────────────────────────────────────────────────────────
\usepackage[hidelinks]{hyperref}

% ─── Matematica ───────────────────────────────────────────────────────────────
\usepackage{amsmath}

% ─── Sezioni compatte ─────────────────────────────────────────────────────────
\usepackage{titlesec}
\titlespacing*{\section}{0pt}{14pt}{6pt}
\titlespacing*{\subsection}{0pt}{10pt}{4pt}

% ─── Header / Footer ──────────────────────────────────────────────────────────
\usepackage{fancyhdr}
\pagestyle{fancy}
\fancyhf{}
\fancyhead[L]{\small HYDRAS --- Report v10}
\fancyhead[R]{\small \today}
\fancyfoot[C]{\thepage}
\renewcommand{\headrulewidth}{0.4pt}

% ─── Codice inline ────────────────────────────────────────────────────────────
\usepackage{listings}
\lstset{basicstyle=\ttfamily\small, breaklines=true}

% ══════════════════════════════════════════════════════════════════════════════
\begin{document}

% ─── Titolo ───────────────────────────────────────────────────────────────────
\begin{center}
    {\LARGE\bfseries HYDRAS --- Report v10}\\[6pt]
    {\large Studio di Ablation della Funzione di Reward}
\end{center}
\vspace{0.3cm}
\hrule
\vspace{0.5cm}

% ══════════════════════════════════════════════════════════════════════════════
\section{Algoritmo PPO}

L'agente \`e addestrato con \textbf{MaskablePPO} (sb3-contrib), estensione di Proximal Policy
Optimization che integra l'action masking per impedire collisioni con la costa.
L'architettura della policy \`e una MLP con due reti separate (attore e critico), ciascuna
composta da due strati nascosti da 256 neuroni con attivazione $\tanh$.
Il training utilizza \textbf{SubprocVecEnv} con 6 worker paralleli (2 worker $\times$ 3 chunk)
avviati con \texttt{start\_method='fork'} per ridurre l'overhead di inizializzazione
e raggiungere un throughput di circa $300\,\text{K}$ step/ora.

% ══════════════════════════════════════════════════════════════════════════════
\section{Dataset --- 132 Sorgenti $\times$ 4 Scenari Vento}

\subsection{Acquisizione e Gestione Dati}

Il dataset \`e composto da simulazioni oceanografiche MIKE21 su un dominio marino di
$3 \times 2.5\,\text{km}$ nella baia di Cecina (costa toscana, UTM32N), con risoluzione
spaziale di $10\,\text{m}$ e passo temporale di $10\,\text{s}$.
Ogni sorgente \`e identificata da un codice univoco (SRC001--SRC132); per ciascuna sorgente
sono disponibili 4 file di concentrazione corrispondenti ai 4 scenari di vento V0--V3.

\subsection{I Quattro Scenari Vento (V0--V3)}

I quattro scenari di vento sono caratterizzati da direzioni e intensit\`a differenti,
generando plume con morfologie diverse. Il modello viene valutato su tutti e 4 gli scenari
per verificarne la robustezza rispetto alle condizioni ambientali.

\subsection{Split Training / Inferenza}

Il dataset \`e suddiviso in modo da garantire una valutazione su sorgenti mai viste
durante il training:
\begin{itemize}
    \item \textbf{Training:} SRC001--SRC106 (80\%, 106 sorgenti)
    \item \textbf{Inferenza:} SRC107--SRC132 (20\%, 26 sorgenti)
\end{itemize}

% ══════════════════════════════════════════════════════════════════════════════
\section{Ambiente di Simulazione}

L'ambiente implementa l'interfaccia \texttt{gymnasium.Env} con spazio delle azioni discreto
a 8 direzioni. Ogni episodio ha una durata massima di 1080 step ($= 3\,\text{h}$ di tempo
simulato). Il dataset \`e suddiviso in 3 chunk temporali (Q1/4, Q1/2, Q3/4) per variare
le condizioni iniziali del campo di concentrazione lungo l'asse temporale della simulazione.

\subsection{Constraints di Spawn e Fallback Automatico}

Lo spawn dell'agente viene campionato uniformemente nelle celle marine del dominio che
soddisfano simultaneamente tre condizioni: (1) assenza di terra secondo la land mask,
(2) distanza dalla costa $\geq 50\,\text{m}$, (3) concentrazione chimica $c \geq 0.5$.
Se nessuna cella soddisfa tutte e tre le condizioni, il threshold di concentrazione
viene abbassato progressivamente fino a trovare una posizione valida (fallback automatico).

% ══════════════════════════════════════════════════════════════════════════════
\section{Training e Architettura}

\subsection{Input della Rete Neurale --- 116 Dimensioni}

L'observation space \`e composto da 116 dimensioni che forniscono all'agente
una rappresentazione completa del campo chimico locale, della propria cinematica
e delle condizioni ambientali (vento e corrente).
Tutte le feature sono normalizzate nell'intervallo $[0,1]$ o $[-1,1]$.

\begin{table}[htbp]
\centering
\small
\begin{tabular}{L{4.5cm} C{1.2cm} L{9cm}}
\toprule
\textbf{Feature} & \textbf{Dim.} & \textbf{Descrizione} \\
\midrule
Concentrazione attuale & 1 & Concentrazione normalizzata nella posizione corrente \\
Memoria concentrazione & 9 & Ultime 9 letture di concentrazione \\
Gradienti spaziali & 8 & Gradiente di concentrazione nelle 8 direzioni \\
Posizione normalizzata & 2 & $x, y$ nel dominio $[0,1]^2$ \\
Velocit\`a agente & 2 & $v_x, v_y$ normalizzati \\
Vento (componenti) & 2 & $u_w, v_w$ normalizzati \\
Corrente (componenti) & 2 & $u_c, v_c$ normalizzati \\
Memoria spostamenti & 18 & Ultimi 9 spostamenti $(dx, dy)$ \\
Sensori direzionali & 8 & Concentrazione nelle 8 direzioni entro 20\,m \\
Massima conc. osservata & 3 & Valore, posizione $(x,y)$ del massimo storico \\
Step senza plume & 1 & Step dall'ultimo contatto con il plume (norm.) \\
Step totali & 1 & Progress episodio normalizzato \\
Distanza dalla costa & 1 & Distanza minima dalla costa normalizzata \\
\bottomrule
\end{tabular}
\caption{Input della rete neurale: 116 dimensioni totali.}
\end{table}

\subsection{Parametri di Training}

I tre modelli dell'ablation study condividono gli stessi iperparametri; l'unica
differenza tra le rispettive configurazioni \`e il parametro \texttt{reward\_mode}.

\begin{table}[htbp]
\centering
\small
\begin{tabular}{L{5cm} C{4cm}}
\toprule
\textbf{Parametro} & \textbf{Valore} \\
\midrule
Algoritmo & MaskablePPO \\
Policy & MlpPolicy (256$\times$256, tanh) \\
Step totali & 6\,000\,000 \\
$n\_steps$ & 4096 \\
Batch size & 128 \\
Epoche & 10 \\
$\gamma$ & 0.99 \\
$\lambda$ (GAE) & 0.95 \\
Clip range & 0.15 \\
$v_f$ coef & 0.3 \\
Ent. coef & 0.01 \\
Max grad norm & 0.5 \\
Target KL & 0.01 \\
\bottomrule
\end{tabular}
\caption{Iperparametri comuni ai tre training dell'ablation study.}
\end{table}

\subsection{Learning Rate Decrescente}

Il learning rate segue uno schedule a gradini decrescenti: una fase iniziale ad alto
learning rate favorisce l'esplorazione, mentre le fasi successive affinano la policy
con aggiornamenti sempre pi\`u conservativi.
Durante il training, la distribuzione degli scenari evolve progressivamente verso
i casi pi\`u difficili tramite un parametro $\alpha \in [0,1]$: per $\alpha = 0$
la distribuzione \`e uniforme, per $\alpha \to 1$ il peso \`e concentrato sugli
scenari con vento pi\`u complesso (V2) e sui chunk temporali centrali.

\begin{table}[htbp]
\centering
\small
\begin{minipage}{0.45\linewidth}
\centering
\begin{tabular}{C{3.5cm} C{2cm}}
\toprule
\textbf{Step} & \textbf{LR} \\
\midrule
0 -- 2\,000\,000 & $3 \times 10^{-4}$ \\
2\,000\,000 -- 4\,000\,000 & $1 \times 10^{-4}$ \\
4\,000\,000 -- 5\,000\,000 & $5 \times 10^{-5}$ \\
5\,000\,000 -- 6\,000\,000 & $1 \times 10^{-5}$ \\
\bottomrule
\end{tabular}
\caption{Schedule del learning rate.}
\end{minipage}
\hfill
\begin{minipage}{0.45\linewidth}
\centering
\begin{tabular}{C{3.5cm} C{2cm}}
\toprule
\textbf{Step} & \textbf{$\alpha$} \\
\midrule
0 -- 2\,000\,000 & 0.0 \\
2\,000\,000 -- 4\,000\,000 & 0.5 \\
4\,000\,000 -- 6\,000\,000 & 0.8 \\
\bottomrule
\end{tabular}
\caption{Curriculum sugli scenari.}
\end{minipage}
\end{table}

\clearpage
% ══════════════════════════════════════════════════════════════════════════════
\section{Studio di Ablation della Funzione di Reward}

\subsection{Motivazione}

\`E stato condotto uno studio di ablation per verificare
l'effetto dei diversi componenti della funzione di reward sulle prestazioni dell'agente.
L'obiettivo \`e determinare quali componenti siano effettivamente necessari e quali
possano invece introdurre segnali di rinforzo fuorvianti.

\subsection{Le Tre Varianti di Reward}

Sono stati definiti tre modi di reward, configurabili tramite il parametro \texttt{reward\_mode}
nel file di configurazione. L'observation space (116 dimensioni) rimane identico in tutti e tre.

\subsubsection*{\texttt{full} --- Reward Completa (modello di riferimento)}

La reward function completa, utilizzata nel modello \texttt{ppo\_20260513}:

\begin{itemize}
    \item Bonus sorgente raggiunta: $+100$ (+ time bonus fino a $+50$)
    \item Penalt\`a uscita dal dominio: $-10$
    \item Reward distanza quadratica con zone multiplier e retreat asymmetry ($\times 3$)
    \item Reward gradiente di concentrazione: $\pm 0.05$
    \item Reward allineamento vento: $\pm 0.05$
    \item Reward allineamento corrente: $\pm 0.05$
    \item Penalt\`a tempo: $-0.1$ per step
    \item Penalt\`a land proximity: fino a $-15$
    \item Stagnation penalty: $-1.5$ (finestra 15 step)
    \item Directional stagnation penalty: $-0.5$
\end{itemize}

\subsubsection*{\texttt{base} --- Reward di Base}

Versione semplificata: rimossi i componenti ad hoc (stagnation, zone multiplier,
retreat asymmetry), distanza lineare al posto di quella quadratica.
Wind e current alignment rimangono attivi.

\begin{itemize}
    \item Bonus sorgente: $+100$ (+ time bonus)
    \item Penalt\`a uscita dominio: $-10$
    \item Reward distanza lineare: $(d_{t-1} - d_t) \times 100 / 3000$
    \item Reward gradiente concentrazione: $\pm 0.05$
    \item Reward allineamento vento: $\pm 0.05$
    \item Reward allineamento corrente: $\pm 0.05$
    \item Penalt\`a tempo: $-0.1$ per step
    \item Penalt\`a land proximity: fino a $-15$
\end{itemize}

\subsubsection*{\texttt{minimal} --- Reward Minima}

Identica a \texttt{base}, con wind e current alignment rimossi dal reward.
Le relative osservazioni rimangono nell'input della rete invariate.

\begin{itemize}
    \item Bonus sorgente: $+100$ (+ time bonus)
    \item Penalt\`a uscita dominio: $-10$
    \item Reward distanza lineare: $(d_{t-1} - d_t) \times 100 / 3000$
    \item Reward gradiente concentrazione: $\pm 0.05$
    \item Penalt\`a tempo: $-0.1$ per step
    \item Penalt\`a land proximity: fino a $-15$
\end{itemize}

\subsection{Setup del Training Sequenziale}

I modelli \texttt{base} e \texttt{minimal} sono stati addestrati
sequenzialmente da zero, ciascuno per 6\,000\,000 step, con gli stessi
iperparametri e lo stesso observation space del modello di riferimento \texttt{full}.
Il training \`e stato avviato automaticamente dallo script \texttt{src/train\_ppo.py}
tramite la funzione \texttt{main\_ablation()}, senza intervento manuale tra i due run.

\begin{table}[htbp]
\centering
\small
\begin{tabular}{L{4.2cm} C{3.8cm} C{4cm}}
\toprule
\textbf{Modello} & \textbf{Config} & \textbf{reward\_mode} \\
\midrule
ppo\_20260513\_154951 & config.yaml & \texttt{full} \\
ppo\_20260516\_095551 & config\_base.yaml & \texttt{base} \\
ppo\_20260516\_143937 & config\_base\_no\_wind.yaml & \texttt{minimal} \\
\bottomrule
\end{tabular}
\caption{Modelli addestrati nell'ablation study con relativa configurazione.}
\end{table}

\subsection{Risultati}

L'inferenza \`e stata condotta sulle 26 sorgenti held-out (SRC107--SRC132),
3 chunk temporali (Q1/4, Q1/2, Q3/4), 4 scenari vento (V0--V3),
5 episodi per scenario --- per un totale di \textbf{1560 episodi} per modello.
La policy \`e stata valutata in modalit\`a deterministica.

\begin{table}[htbp]
\centering
\renewcommand{\arraystretch}{1.3}
\small
\newcommand{\dpos}[1]{{\scriptsize\textcolor{green!50!black}{(+#1\%)}}}
\newcommand{\dneg}[1]{{\scriptsize\textcolor{red!70!black}{(-#1\%)}}}
\newcommand{\dzero}{{\scriptsize\textcolor{gray}{(0\%)}}}
\begin{tabular}{L{3.5cm} C{2.4cm} C{3.2cm} C{3.2cm}}
\toprule
\rowcolor{headerbg}
\textcolor{white}{\textbf{Metrica}} &
\textcolor{white}{\textbf{base}} &
\textcolor{white}{\textbf{full}} &
\textcolor{white}{\textbf{minimal}} \\
\midrule
\rowcolor{rowalt}
\textbf{SR globale} & 88{,}3\% & 95{,}6\% \dpos{7{,}3} & \textbf{96{,}9\%} \dpos{1{,}3} \\
Dist. iniziale media & 799{,}9\,m & 802{,}1\,m & 801{,}1\,m \\
\rowcolor{rowalt}
Step medi (successo) & 113{,}8 ($\sim$19{,}0\,min) & 129{,}6 ($\sim$21{,}6\,min) & 129{,}2 ($\sim$21{,}5\,min) \\
\midrule
V0 & 99{,}5\% & 98{,}0\% \dneg{1{,}5} & 100{,}0\% \dpos{2{,}0} \\
\rowcolor{rowalt}
V1 & 95{,}4\% & 99{,}0\% \dpos{3{,}6} & 99{,}0\% \dzero \\
V2 & 67{,}9\% & 88{,}5\% \dpos{20{,}6} & 88{,}7\% \dpos{0{,}2} \\
\rowcolor{rowalt}
V3 & 90{,}3\% & 97{,}0\% \dpos{6{,}7} & 99{,}7\% \dpos{2{,}7} \\
\midrule
Q1/4 & 100{,}0\% & 100{,}0\% \dzero & 100{,}0\% \dzero \\
\rowcolor{rowalt}
Q1/2 & 74{,}4\% & 93{,}7\% \dpos{19{,}3} & 93{,}5\% \dneg{0{,}2} \\
Q3/4 & 90{,}4\% & 93{,}0\% \dpos{2{,}6} & 97{,}1\% \dpos{4{,}1} \\
\bottomrule
\end{tabular}
\caption{Confronto delle prestazioni dei tre modelli su 1560 episodi ciascuno
         (26 sorgenti held-out, 4 venti, 3 chunk, 5 episodi).
         I delta sono calcolati rispetto alla colonna precedente.}
\end{table}

\subsection{Analisi dei Risultati}

I grafici seguenti mostrano le distribuzioni di distanza iniziale, success rate, tempo al successo e distanza residua nel tempo per i modelli \texttt{base} e \texttt{minimal}, calcolati sui 1560 episodi di inferenza.

\subsubsection*{Modello \texttt{base}}

\begin{figure}[htbp]
  \centering
  \begin{minipage}[t]{0.48\textwidth}
    \centering
    \includegraphics[width=\linewidth]{../../evaluations/evaluations_ablation_base/analysis/plot_initial_dist.png}
    \caption{Distribuzione delle distanze iniziali --- \texttt{base}.}
  \end{minipage}\hfill
  \begin{minipage}[t]{0.48\textwidth}
    \centering
    \includegraphics[width=\linewidth]{../../evaluations/evaluations_ablation_base/analysis/plot_sr_analysis.png}
    \caption{Success rate per scenario --- \texttt{base}.}
  \end{minipage}
\end{figure}

\begin{figure}[htbp]
  \centering
  \begin{minipage}[t]{0.48\textwidth}
    \centering
    \includegraphics[width=\linewidth]{../../evaluations/evaluations_ablation_base/analysis/plot_success_time_dist.png}
    \caption{Distribuzione del tempo al successo --- \texttt{base}.}
  \end{minipage}\hfill
  \begin{minipage}[t]{0.48\textwidth}
    \centering
    \includegraphics[width=\linewidth]{../../evaluations/evaluations_ablation_base/analysis/plot_distance_over_time.png}
    \caption{Distanza media dalla sorgente nel tempo --- \texttt{base}.}
  \end{minipage}
\end{figure}

\subsubsection*{Modello \texttt{minimal}}

\begin{figure}[htbp]
  \centering
  \begin{minipage}[t]{0.48\textwidth}
    \centering
    \includegraphics[width=\linewidth]{../../evaluations/evaluations_ablation_base_no_wind_reward/analysis/plot_initial_dist.png}
    \caption{Distribuzione delle distanze iniziali --- \texttt{minimal}.}
  \end{minipage}\hfill
  \begin{minipage}[t]{0.48\textwidth}
    \centering
    \includegraphics[width=\linewidth]{../../evaluations/evaluations_ablation_base_no_wind_reward/analysis/plot_sr_analysis.png}
    \caption{Success rate per scenario --- \texttt{minimal}.}
  \end{minipage}
\end{figure}

\begin{figure}[htbp]
  \centering
  \begin{minipage}[t]{0.48\textwidth}
    \centering
    \includegraphics[width=\linewidth]{../../evaluations/evaluations_ablation_base_no_wind_reward/analysis/plot_success_time_dist.png}
    \caption{Distribuzione del tempo al successo --- \texttt{minimal}.}
  \end{minipage}\hfill
  \begin{minipage}[t]{0.48\textwidth}
    \centering
    \includegraphics[width=\linewidth]{../../evaluations/evaluations_ablation_base_no_wind_reward/analysis/plot_distance_over_time.png}
    \caption{Distanza media dalla sorgente nel tempo --- \texttt{minimal}.}
  \end{minipage}
\end{figure}

\subsection{Conclusioni}

\subsubsection*{Reward minima vs.\ reward completa}

Il modello \texttt{minimal} ottiene un success rate del \textbf{97\%},
superando sia il modello \texttt{full} (96\%) che \texttt{base} (88\%),
con configurazione di training identica e stesso observation space.
Questo risultato indica che i componenti aggiuntivi della reward \texttt{full}
(stagnation penalty, zone multiplier, retreat asymmetry) non producono un guadagno
misurabile rispetto a una reward lineare pi\`u semplice.

\subsubsection*{Effetto del wind/current alignment}

Il delta pi\`u significativo \`e tra \texttt{base} e \texttt{minimal}:
$-9\%$ di success rate globale, con un crollo di circa 20 punti percentuali
su Q1/2 e su V2.
I reward di allineamento vento e corrente ($\pm 0.05$ ciascuno) assumono che la sorgente
si trovi sempre controvento/controcorrente rispetto all'agente. Questa ipotesi non \`e
valida in un ambiente costiero complesso, dove riflessioni, correnti locali e
morfologia della costa rendono il percorso ottimale non riducibile alla direzione del flusso.
Il segnale risultante \`e in conflitto con il gradiente chimico in una frazione rilevante
degli scenari, degradando la politica appresa.

Si osserva inoltre che i componenti aggiuntivi della reward \texttt{full}
(stagnation, zone multiplier, retreat asymmetry) sembrano svolgere in parte il ruolo di
compensare il segnale negativo introdotto dal wind/current alignment, portando il
risultato finale a un livello comparabile con \texttt{minimal} nonostante
la presenza dello stesso segnale fuorviante.

\subsubsection*{Ruolo del gradiente chimico}

Il gradiente di concentrazione ($\pm 0.05$) \`e l'unico segnale direzionale
nella reward \texttt{minimal}. Essendo calcolato direttamente sul campo
di concentrazione simulato, incorpora implicitamente l'effetto di vento e corrente ---
poich\'e il plume si disperde seguendo questi flussi. L'agente non necessita di ricevere
l'informazione sul vento nel reward se pu\`o gi\`a leggerla nella struttura spaziale
del campo chimico.

\subsubsection*{Osservazioni sulla reward engineering}

I risultati suggeriscono un principio generale per questo tipo di problema:
\textbf{un observation space ricco rende il reward shaping complesso ridondante}.
I componenti aggiuntivi del modello \texttt{full} erano stati introdotti iterativamente
per correggere failure mode osservati con versioni precedenti dell'observation space.
Con 116 dimensioni di osservazione --- in particolare il gradiente chimico direzionale,
la memoria degli spostamenti e la storia della concentrazione --- l'agente dispone di
informazioni sufficienti per navigare efficacemente senza spinte esplicite aggiuntive
nel reward.

\end{document}
